\documentclass[11pt]{article}
\usepackage{icmlko}
\begin{document}

\title{순위를 매기던 자가 다섯 중 제일 약했다 --- 판 차 $83$ 건 중 $60$ 건이 표본 잡음 안이다}
\author{월드모델 실험실}
\date{노트 611 · 2026년 8월 5일}
\maketitle

\begin{abstract}
이 판의 상시 규약은 ``\textbf{순위는 유보 표본의 판 $\rho$}''다. 노트
$609$ 가 판의 \textbf{순수 표본 SD 를 $0.0023$} 으로 쟀다 --- 씨앗을
늘려도 안 줄어드는 성분이다. 원장을 그 자로 훑으니

\begin{center}
판 차를 적은 \textbf{$83$ 건 중 $60$ 건}$(72\%)$이 \textbf{$2$ SD
$=0.0045$ 아래}
\end{center}

\textbf{순위를 매기던 자가 대부분의 결정을 못 가른다.} 그런데 이 판은
자가 다섯이다. 채택을 실제로 정한 \textbf{집 밖 두 짝}을 같은 자로 쟀다.

\begin{center}
\begin{tabular}{lrrr}
\hline
자 & 효과 & 표본 SD & $t$ \\
\hline
판($3{,}369$ 행) & $+0.0018$ & $\mathbf{0.0023}$ & $\mathbf{0.8}$ \\
KR 만화($1{,}716$) & $+0.0390$ & $0.0128$ & $3.0$ \\
비게임 앱($1{,}600$) & $+0.2187$ & $0.0122$ & $\mathbf{17.9}$ \\
\hline
\end{tabular}
\end{center}

\textbf{판은 잡음이 제일 작지만 효과도 제일 작아서 신호 대 잡음이 제일
나쁘다.} 다섯 자를 둔 설계가 이 문제를 미리 막고 있었다 --- \emph{그것을
쟀던 적은 없다.}
\end{abstract}

\begin{figure}[t]
\centering
\includegraphics[width=\textwidth]{figs/f1.pdf}
\caption{\textbf{가} --- 판은 잡음도 효과도 제일 작다.
\textbf{나} --- 판만 $0$ 을 문다.
\textbf{다} --- 원장이 적은 판 차 대부분이 잡음 안이다.}
\end{figure}

\section{잡음은 유보의 성질이고 자마다 다르다}

노트 $610$ 이 판에서 두 가지 팔(같은 적합에서 열만 가림 $\cdot$ 두 번
적합)의 \textbf{순수 표본 SD 가 $0.00227$ 과 $0.00229$ 로 사실상 같음}을
봤다. 표본 성분은 \emph{팔}이 아니라 \emph{유보}의 성질이다. 그래서 판에서
$0.0023$ 을 상수로 쓸 수 있다.

집 밖 짝은 다르다 --- $0.0128 \cdot 0.0122$ 로 \textbf{판의 다섯 배}다.
한 도메인이고 행이 절반이니 당연하다. \textbf{그런데 효과가 스무 배
$\cdot$ 백 배 크다.}

\begin{quote}
\textbf{잡음이 작은 자가 좋은 자가 아니다.} 판은 열한 도메인을 평균해
잡음을 눌렀지만, 같은 평균이 \emph{효과도} 눌렀다 --- 노트 $610$ 이 봤듯
채택의 판 이득은 고친 넷 $-0.0009$ 와 이웃 $+0.0031$ 이 상쇄된
\textbf{잔차}였다.
\end{quote}

\section{원장 $83$ 건}

원장에서 ``채택''이 들어간 항목 중 판 차를 적은 것이 $83$ 건이다(글자
훑기라 기각 $\cdot$ 사전 등록이 섞여 있다 --- 정확한 수가 아니라
\emph{분포}로 읽는다).

\begin{center}
\begin{tabular}{lr}
\hline
$|$판 차$|$ & 건수 \\
\hline
$\ge 2$ SD $(0.0045)$ & $23$ \\
$< 2$ SD & $\mathbf{60}$ \\
$< 1$ SD $(0.0023)$ & $38$ \\
\hline
\end{tabular}
\end{center}

$2$ SD 위 $23$ 건은 대체로 \textbf{누수 $\cdot$ 오염 $\cdot$ 수집}이다
(웹툰 검색 마스크 $-0.48$ $\cdot$ 펀딩 표집 $-0.060$ $\cdot$ 시계열 수집
$-0.028$ $\cdot$ 날짜 구멍 $+0.019$). \textbf{모형 손잡이는 거의 다 아래
칸에 있다} --- 깊이 $6{\to}12$ $-0.0012$, 계절 축 $+0.0041$, \texttt{grp}
축 $+0.0037$.

\begin{quote}
\textbf{자료를 고치면 판이 크게 움직이고, 모형을 고치면 잡음 안에서
움직인다.} 이 판이 두 해 동안 자료 쪽에서 훨씬 크게 벌었다는 것을
\emph{잡음의 자}로 다시 확인한 셈이다.
\end{quote}

\section{규약은 산다 --- 다만 사정거리가 있다}

\begin{quote}
\textbf{판 차가 $0.0045$ 아래면 판은 순위를 못 정한다.} 그 구간에서는
집 밖 두 짝이 정한다 --- 노트 $586$ 채택이 정확히 그 예다. 판은
$t{=}0.8$ 로 $0$ 을 물고, KR 만화 $t{=}3.0$ $\cdot$ 비게임 앱
$t{=}17.9$ 로 \textbf{둘 다 $0$ 밖}이다. \textbf{채택은 흔들리지 않는다.}
\end{quote}

이것이 노트 $426$ 의 거울 조항(집 밖 두 짝이 같은 부호 $\cdot$ 둘 다
구간이 $0$ 을 안 묾)이 왜 옳았는지를 설명한다. 그 조항은 \emph{전이}를
지키려고 세웠는데, 실제로는 \textbf{판이 못 가르는 구간에서 유일하게
남는 자}였다.

\section*{조항}

\begin{quote}
\textbf{자의 잡음을 재고 나서 그 자로 순위를 매긴다.} ``판 $\rho$ 로
순위''는 판 차가 $2$ SD $(0.0045)$ 위일 때만 유효하고, 원장이 적은
$83$ 건 중 $60$ 건이 그 아래다. \textbf{그 아래에서는 집 밖 짝이 정한다.}
잡음이 제일 작은 자가 제일 좋은 자라는 직관은 틀렸다 --- \textbf{효과도
같이 눌리기 때문이다.}
\end{quote}

\section*{남은 것}

\begin{itemize}
\item \textbf{날짜 통제 판 $\cdot$ 무리 안 판은 안 쟀다} --- 다섯 자 중
셋만 쟀다. 둘은 판과 같은 행을 쓰므로 SD 가 비슷하겠지만 \emph{안 쟀다.}
\item \textbf{$60$ 건을 하나씩 다시 보지 않았다} --- 판만으로 채택한
것이 몇 건인지, 그중 집 밖 짝이 반대 부호였던 것이 있는지 안 봤다.
\textbf{여기가 다음이다.}
\item \textbf{원장 훑기가 글자 기반}이다. 기각과 사전 등록이 섞여 있어
$83 \cdot 60$ 은 정확한 수가 아니다. 원장에 판 차를 \emph{필드로}
적었으면 셀 수 있었다.
\end{itemize}

\end{document}
